
\section{Conclusion}
\label{sec:conclusion}

We have presented \thecalculus, a calculus of context-free session
types with borrowing.  The central operation is session-type
splitting: a borrow consumes a prefix of a protocol and leaves a
residual suffix, but for CFSTs this operation must be understood
modulo the equational theory of sequential composition, choices,
skipping behavior, and recursion.  Borrowing therefore lifts the
programming convenience of temporary access from regular session types
to the richer protocol structure of CFSTs, while preserving the
linear view that using an endpoint advances its protocol state.

The calculus separates local from remote borrowing.  Local borrows
stay within a process and can be sequenced by ordered typing and
evaluation alone; remote borrows may cross process boundaries and
therefore require an explicit handoff between the borrowed prefix and
the residual suffix.  This distinction is reflected throughout the
development: in the declarative type system, in the direct operational
semantics based on binder groups, and in the translation to a lower
level calculus with primitive channels and synchronization cells.  The
algorithmic type system reduces the splitting obligations to
constraints over CFST equivalence, for which our implementation
reuses the FreeST solver.

The resulting picture is that borrowing is not merely syntactic
sugar for passing endpoints around.  In the presence of context-free
session types, it is a disciplined way to compute and schedule pieces
of a communication protocol.  Several directions remain open,
including stronger inference for split constraints, richer programming
interfaces that avoid undecidable subtyping, and integration with
larger CFST implementations.


%%% Local Variables:
%%% mode: latex
%%% TeX-master: "../popl2027.tex"
%%% End:
